\documentclass[10pt]{article}
\usepackage[margin=0.7in]{geometry}
\usepackage{booktabs}
\usepackage{longtable}
\usepackage{lmodern}
\title{HumanEval TEST-set typicality metrics\\ \large gemma-4-31b \& qwen-3.5-9b $\times$ base/s1/s2/s3/s4/s7/s13}
\author{}
\date{2026-06-18}
\begin{document}
\maketitle

One table per metric (gen ROC, val ROC, val acc, Pearson). Each shows the \textbf{raw} and
\textbf{tc} (typicality-corrected) gen variants, for all systems on both
\texttt{humaneval-v2.1correct-upper} and \texttt{-multi}. Values $\times 100$, 1 decimal;
means over the 82 test problems. Eval = (typicality mode)/(base-typ vs no-base).

\medskip\noindent\textit{Notes:} (i) \texttt{val ROC} and \texttt{val acc} are validator-based, so
their raw and tc columns are identical by construction (shown for uniformity). (ii) \texttt{raw}
gen scores do not depend on the eval mode, so the raw column repeats within a (model, setting,
dataset) block. (iii) gemma base test eval was not in the harvested dirs (qwen base is present).

{\small
\begin{longtable}{llll rr}
\caption{HumanEval TEST gen ROC ($\times100$), raw vs tc gen variant. Systems: gemma-4-31b \& qwen-3.5-9b $\times$ base/s1/s2/s3/s4/s7/s13, both datasets.}\label{tab:he_test_gen_roc}\\
\toprule
Model & Data & Setting & Eval & raw & tc \\
\midrule
\endfirsthead
\multicolumn{6}{l}{\itshape (continued)}\\
\toprule
Model & Data & Setting & Eval & raw & tc \\
\midrule
\endhead
\bottomrule
\endfoot
gemma-4-31b & upper & s1 & self/base-typ & 81.7 & 82.9 \\
gemma-4-31b & upper & s1 & self/no-base & 81.7 & 80.6 \\
gemma-4-31b & upper & s1 & neg/base-typ & 81.7 & 81.6 \\
gemma-4-31b & upper & s1 & neg/no-base & 81.7 & 76.5 \\
gemma-4-31b & upper & s2 & self/base-typ & 90.7 & 86.9 \\
gemma-4-31b & upper & s2 & self/no-base & 90.7 & 85.2 \\
gemma-4-31b & upper & s2 & neg/base-typ & 90.7 & 88.2 \\
gemma-4-31b & upper & s2 & neg/no-base & 90.7 & 86.3 \\
gemma-4-31b & upper & s3 & self/base-typ & 91.5 & 87.5 \\
gemma-4-31b & upper & s3 & self/no-base & 91.5 & 86.4 \\
gemma-4-31b & upper & s3 & neg/base-typ & 91.5 & 87.9 \\
gemma-4-31b & upper & s3 & neg/no-base & 91.5 & 82.3 \\
gemma-4-31b & upper & s4 & self/base-typ & 85.9 & 90.7 \\
gemma-4-31b & upper & s4 & self/no-base & 85.9 & 84.3 \\
gemma-4-31b & upper & s7 & neg/base-typ & 83.6 & 94.8 \\
gemma-4-31b & upper & s7 & neg/no-base & 83.6 & 79.1 \\
gemma-4-31b & upper & s13 & self/base-typ & 80.1 & 83.4 \\
gemma-4-31b & upper & s13 & self/no-base & 80.1 & 79.9 \\
gemma-4-31b & upper & s13 & neg/base-typ & 80.1 & 84.4 \\
gemma-4-31b & upper & s13 & neg/no-base & 80.1 & 75.2 \\
\midrule
gemma-4-31b & multi & s1 & self/base-typ & 66.4 & 87.9 \\
gemma-4-31b & multi & s1 & self/no-base & 66.4 & 81.6 \\
gemma-4-31b & multi & s1 & neg/base-typ & 66.4 & 93.3 \\
gemma-4-31b & multi & s1 & neg/no-base & 66.4 & 72.1 \\
gemma-4-31b & multi & s2 & self/base-typ & 85.6 & 90.7 \\
gemma-4-31b & multi & s2 & self/no-base & 85.6 & 85.6 \\
gemma-4-31b & multi & s2 & neg/base-typ & 85.6 & 94.2 \\
gemma-4-31b & multi & s2 & neg/no-base & 85.6 & 82.4 \\
gemma-4-31b & multi & s3 & self/base-typ & 87.2 & 90.9 \\
gemma-4-31b & multi & s3 & self/no-base & 87.2 & 89.0 \\
gemma-4-31b & multi & s3 & neg/base-typ & 87.2 & 94.0 \\
gemma-4-31b & multi & s3 & neg/no-base & 87.2 & 86.0 \\
gemma-4-31b & multi & s4 & self/base-typ & 74.2 & 90.1 \\
gemma-4-31b & multi & s4 & self/no-base & 74.2 & 84.0 \\
gemma-4-31b & multi & s7 & neg/base-typ & 68.1 & 92.0 \\
gemma-4-31b & multi & s7 & neg/no-base & 68.1 & 80.2 \\
gemma-4-31b & multi & s13 & self/base-typ & 64.2 & 87.3 \\
gemma-4-31b & multi & s13 & self/no-base & 64.2 & 80.1 \\
gemma-4-31b & multi & s13 & neg/base-typ & 64.2 & 93.3 \\
gemma-4-31b & multi & s13 & neg/no-base & 64.2 & 72.4 \\
\midrule
qwen-3.5-9b & upper & s1 & self/base-typ & 76.2 & 87.9 \\
qwen-3.5-9b & upper & s1 & self/no-base & 76.2 & 84.1 \\
qwen-3.5-9b & upper & s1 & neg/base-typ & 76.2 & 94.1 \\
qwen-3.5-9b & upper & s1 & neg/no-base & 76.2 & 73.9 \\
qwen-3.5-9b & upper & s2 & self/base-typ & 83.4 & 87.3 \\
qwen-3.5-9b & upper & s2 & self/no-base & 83.4 & 82.2 \\
qwen-3.5-9b & upper & s2 & neg/base-typ & 83.4 & 90.2 \\
qwen-3.5-9b & upper & s2 & neg/no-base & 83.4 & 79.3 \\
qwen-3.5-9b & upper & s3 & self/base-typ & 81.3 & 86.8 \\
qwen-3.5-9b & upper & s3 & self/no-base & 81.3 & 83.0 \\
qwen-3.5-9b & upper & s3 & neg/base-typ & 81.3 & 92.6 \\
qwen-3.5-9b & upper & s3 & neg/no-base & 81.3 & 80.3 \\
qwen-3.5-9b & upper & s4 & self/base-typ & 69.0 & 83.2 \\
qwen-3.5-9b & upper & s4 & self/no-base & 69.0 & 78.6 \\
qwen-3.5-9b & upper & s7 & neg/base-typ & 65.8 & 88.9 \\
qwen-3.5-9b & upper & s7 & neg/no-base & 65.8 & 77.5 \\
qwen-3.5-9b & upper & s13 & self/base-typ & 74.1 & 86.0 \\
qwen-3.5-9b & upper & s13 & self/no-base & 74.1 & 81.9 \\
qwen-3.5-9b & upper & s13 & neg/base-typ & 74.1 & 92.6 \\
qwen-3.5-9b & upper & s13 & neg/no-base & 74.1 & 72.7 \\
\midrule
qwen-3.5-9b & multi & s1 & self/base-typ & 59.0 & 89.4 \\
qwen-3.5-9b & multi & s1 & self/no-base & 59.0 & 83.3 \\
qwen-3.5-9b & multi & s1 & neg/base-typ & 59.0 & 94.5 \\
qwen-3.5-9b & multi & s1 & neg/no-base & 59.0 & 66.1 \\
qwen-3.5-9b & multi & s2 & self/base-typ & 86.7 & 93.1 \\
qwen-3.5-9b & multi & s2 & self/no-base & 86.7 & 90.2 \\
qwen-3.5-9b & multi & s2 & neg/base-typ & 86.7 & 96.5 \\
qwen-3.5-9b & multi & s2 & neg/no-base & 86.7 & 80.3 \\
qwen-3.5-9b & multi & s3 & self/base-typ & 80.0 & 92.3 \\
qwen-3.5-9b & multi & s3 & self/no-base & 80.0 & 88.5 \\
qwen-3.5-9b & multi & s3 & neg/base-typ & 80.0 & 96.7 \\
qwen-3.5-9b & multi & s3 & neg/no-base & 80.0 & 75.6 \\
qwen-3.5-9b & multi & s4 & self/base-typ & 58.3 & 89.8 \\
qwen-3.5-9b & multi & s4 & self/no-base & 58.3 & 82.9 \\
qwen-3.5-9b & multi & s7 & neg/base-typ & 45.0 & 87.9 \\
qwen-3.5-9b & multi & s7 & neg/no-base & 45.0 & 61.8 \\
qwen-3.5-9b & multi & s13 & self/base-typ & 56.6 & 88.8 \\
qwen-3.5-9b & multi & s13 & self/no-base & 56.6 & 82.9 \\
qwen-3.5-9b & multi & s13 & neg/base-typ & 56.6 & 93.9 \\
qwen-3.5-9b & multi & s13 & neg/no-base & 56.6 & 52.8 \\
\midrule
qwen-3.5-9b (base) & upper & base & self/no-base & 53.6 & 69.1 \\
qwen-3.5-9b (base) & upper & base & neg/no-base & 53.6 & 65.7 \\
\midrule
qwen-3.5-9b (base) & multi & base & self/no-base & 31.3 & 69.0 \\
qwen-3.5-9b (base) & multi & base & neg/no-base & 31.2 & 59.9 \\
\bottomrule
\end{longtable}
}

{\small
\begin{longtable}{llll rr}
\caption{HumanEval TEST val ROC ($\times100$), raw vs tc gen variant. Systems: gemma-4-31b \& qwen-3.5-9b $\times$ base/s1/s2/s3/s4/s7/s13, both datasets.}\label{tab:he_test_val_roc}\\
\toprule
Model & Data & Setting & Eval & raw & tc \\
\midrule
\endfirsthead
\multicolumn{6}{l}{\itshape (continued)}\\
\toprule
Model & Data & Setting & Eval & raw & tc \\
\midrule
\endhead
\bottomrule
\endfoot
gemma-4-31b & upper & s1 & self/base-typ & 93.9 & 93.9 \\
gemma-4-31b & upper & s1 & self/no-base & 93.9 & 93.9 \\
gemma-4-31b & upper & s1 & neg/base-typ & 93.9 & 93.9 \\
gemma-4-31b & upper & s1 & neg/no-base & 93.9 & 93.9 \\
gemma-4-31b & upper & s2 & self/base-typ & 91.4 & 91.4 \\
gemma-4-31b & upper & s2 & self/no-base & 91.4 & 91.4 \\
gemma-4-31b & upper & s2 & neg/base-typ & 91.4 & 91.4 \\
gemma-4-31b & upper & s2 & neg/no-base & 91.4 & 91.4 \\
gemma-4-31b & upper & s3 & self/base-typ & 94.5 & 94.5 \\
gemma-4-31b & upper & s3 & self/no-base & 94.5 & 94.5 \\
gemma-4-31b & upper & s3 & neg/base-typ & 94.5 & 94.5 \\
gemma-4-31b & upper & s3 & neg/no-base & 94.5 & 94.5 \\
gemma-4-31b & upper & s4 & self/base-typ & 93.1 & 93.1 \\
gemma-4-31b & upper & s4 & self/no-base & 93.1 & 93.1 \\
gemma-4-31b & upper & s7 & neg/base-typ & 94.0 & 94.0 \\
gemma-4-31b & upper & s7 & neg/no-base & 94.0 & 94.0 \\
gemma-4-31b & upper & s13 & self/base-typ & 93.9 & 93.9 \\
gemma-4-31b & upper & s13 & self/no-base & 93.9 & 93.9 \\
gemma-4-31b & upper & s13 & neg/base-typ & 93.9 & 93.9 \\
gemma-4-31b & upper & s13 & neg/no-base & 93.9 & 93.9 \\
\midrule
gemma-4-31b & multi & s1 & self/base-typ & 96.1 & 96.1 \\
gemma-4-31b & multi & s1 & self/no-base & 96.1 & 96.1 \\
gemma-4-31b & multi & s1 & neg/base-typ & 96.1 & 96.1 \\
gemma-4-31b & multi & s1 & neg/no-base & 96.1 & 96.1 \\
gemma-4-31b & multi & s2 & self/base-typ & 90.0 & 90.0 \\
gemma-4-31b & multi & s2 & self/no-base & 90.0 & 90.0 \\
gemma-4-31b & multi & s2 & neg/base-typ & 90.0 & 90.0 \\
gemma-4-31b & multi & s2 & neg/no-base & 90.0 & 90.0 \\
gemma-4-31b & multi & s3 & self/base-typ & 93.8 & 93.8 \\
gemma-4-31b & multi & s3 & self/no-base & 93.8 & 93.8 \\
gemma-4-31b & multi & s3 & neg/base-typ & 93.8 & 93.8 \\
gemma-4-31b & multi & s3 & neg/no-base & 93.8 & 93.8 \\
gemma-4-31b & multi & s4 & self/base-typ & 94.4 & 94.4 \\
gemma-4-31b & multi & s4 & self/no-base & 94.4 & 94.4 \\
gemma-4-31b & multi & s7 & neg/base-typ & 94.3 & 94.3 \\
gemma-4-31b & multi & s7 & neg/no-base & 94.3 & 94.3 \\
gemma-4-31b & multi & s13 & self/base-typ & 93.8 & 93.8 \\
gemma-4-31b & multi & s13 & self/no-base & 93.8 & 93.8 \\
gemma-4-31b & multi & s13 & neg/base-typ & 93.8 & 93.8 \\
gemma-4-31b & multi & s13 & neg/no-base & 93.8 & 93.8 \\
\midrule
qwen-3.5-9b & upper & s1 & self/base-typ & 86.5 & 86.5 \\
qwen-3.5-9b & upper & s1 & self/no-base & 86.4 & 86.4 \\
qwen-3.5-9b & upper & s1 & neg/base-typ & 86.5 & 86.5 \\
qwen-3.5-9b & upper & s1 & neg/no-base & 86.4 & 86.4 \\
qwen-3.5-9b & upper & s2 & self/base-typ & 83.3 & 83.3 \\
qwen-3.5-9b & upper & s2 & self/no-base & 83.3 & 83.3 \\
qwen-3.5-9b & upper & s2 & neg/base-typ & 83.2 & 83.2 \\
qwen-3.5-9b & upper & s2 & neg/no-base & 83.3 & 83.3 \\
qwen-3.5-9b & upper & s3 & self/base-typ & 87.6 & 87.6 \\
qwen-3.5-9b & upper & s3 & self/no-base & 87.6 & 87.6 \\
qwen-3.5-9b & upper & s3 & neg/base-typ & 87.6 & 87.6 \\
qwen-3.5-9b & upper & s3 & neg/no-base & 87.6 & 87.6 \\
qwen-3.5-9b & upper & s4 & self/base-typ & 86.3 & 86.3 \\
qwen-3.5-9b & upper & s4 & self/no-base & 86.3 & 86.3 \\
qwen-3.5-9b & upper & s7 & neg/base-typ & 86.1 & 86.1 \\
qwen-3.5-9b & upper & s7 & neg/no-base & 86.0 & 86.0 \\
qwen-3.5-9b & upper & s13 & self/base-typ & 84.6 & 84.6 \\
qwen-3.5-9b & upper & s13 & self/no-base & 84.7 & 84.7 \\
qwen-3.5-9b & upper & s13 & neg/base-typ & 84.6 & 84.6 \\
qwen-3.5-9b & upper & s13 & neg/no-base & 84.6 & 84.6 \\
\midrule
qwen-3.5-9b & multi & s1 & self/base-typ & 90.4 & 90.4 \\
qwen-3.5-9b & multi & s1 & self/no-base & 90.4 & 90.4 \\
qwen-3.5-9b & multi & s1 & neg/base-typ & 90.4 & 90.4 \\
qwen-3.5-9b & multi & s1 & neg/no-base & 90.4 & 90.4 \\
qwen-3.5-9b & multi & s2 & self/base-typ & 84.8 & 84.8 \\
qwen-3.5-9b & multi & s2 & self/no-base & 84.8 & 84.8 \\
qwen-3.5-9b & multi & s2 & neg/base-typ & 84.8 & 84.8 \\
qwen-3.5-9b & multi & s2 & neg/no-base & 84.8 & 84.8 \\
qwen-3.5-9b & multi & s3 & self/base-typ & 89.6 & 89.6 \\
qwen-3.5-9b & multi & s3 & self/no-base & 89.5 & 89.5 \\
qwen-3.5-9b & multi & s3 & neg/base-typ & 89.5 & 89.5 \\
qwen-3.5-9b & multi & s3 & neg/no-base & 89.5 & 89.5 \\
qwen-3.5-9b & multi & s4 & self/base-typ & 89.1 & 89.1 \\
qwen-3.5-9b & multi & s4 & self/no-base & 89.1 & 89.1 \\
qwen-3.5-9b & multi & s7 & neg/base-typ & 88.5 & 88.5 \\
qwen-3.5-9b & multi & s7 & neg/no-base & 88.5 & 88.5 \\
qwen-3.5-9b & multi & s13 & self/base-typ & 88.0 & 88.0 \\
qwen-3.5-9b & multi & s13 & self/no-base & 88.0 & 88.0 \\
qwen-3.5-9b & multi & s13 & neg/base-typ & 88.0 & 88.0 \\
qwen-3.5-9b & multi & s13 & neg/no-base & 88.0 & 88.0 \\
\midrule
qwen-3.5-9b (base) & upper & base & self/no-base & 81.7 & 81.7 \\
qwen-3.5-9b (base) & upper & base & neg/no-base & 81.7 & 81.7 \\
\midrule
qwen-3.5-9b (base) & multi & base & self/no-base & 83.4 & 83.4 \\
qwen-3.5-9b (base) & multi & base & neg/no-base & 83.4 & 83.4 \\
\bottomrule
\end{longtable}
}

{\small
\begin{longtable}{llll rr}
\caption{HumanEval TEST val acc ($\times100$), raw vs tc gen variant. Systems: gemma-4-31b \& qwen-3.5-9b $\times$ base/s1/s2/s3/s4/s7/s13, both datasets.}\label{tab:he_test_val_acc}\\
\toprule
Model & Data & Setting & Eval & raw & tc \\
\midrule
\endfirsthead
\multicolumn{6}{l}{\itshape (continued)}\\
\toprule
Model & Data & Setting & Eval & raw & tc \\
\midrule
\endhead
\bottomrule
\endfoot
gemma-4-31b & upper & s1 & self/base-typ & 79.6 & 79.6 \\
gemma-4-31b & upper & s1 & self/no-base & 79.6 & 79.6 \\
gemma-4-31b & upper & s1 & neg/base-typ & 79.6 & 79.6 \\
gemma-4-31b & upper & s1 & neg/no-base & 79.6 & 79.6 \\
gemma-4-31b & upper & s2 & self/base-typ & 76.8 & 76.8 \\
gemma-4-31b & upper & s2 & self/no-base & 76.8 & 76.8 \\
gemma-4-31b & upper & s2 & neg/base-typ & 76.8 & 76.8 \\
gemma-4-31b & upper & s2 & neg/no-base & 76.8 & 76.8 \\
gemma-4-31b & upper & s3 & self/base-typ & 87.8 & 87.8 \\
gemma-4-31b & upper & s3 & self/no-base & 87.8 & 87.8 \\
gemma-4-31b & upper & s3 & neg/base-typ & 87.8 & 87.8 \\
gemma-4-31b & upper & s3 & neg/no-base & 87.8 & 87.8 \\
gemma-4-31b & upper & s4 & self/base-typ & 84.1 & 84.1 \\
gemma-4-31b & upper & s4 & self/no-base & 84.1 & 84.1 \\
gemma-4-31b & upper & s7 & neg/base-typ & 88.3 & 88.3 \\
gemma-4-31b & upper & s7 & neg/no-base & 88.3 & 88.3 \\
gemma-4-31b & upper & s13 & self/base-typ & 87.9 & 87.9 \\
gemma-4-31b & upper & s13 & self/no-base & 87.9 & 87.9 \\
gemma-4-31b & upper & s13 & neg/base-typ & 87.9 & 87.9 \\
gemma-4-31b & upper & s13 & neg/no-base & 87.9 & 87.9 \\
\midrule
gemma-4-31b & multi & s1 & self/base-typ & 87.2 & 87.2 \\
gemma-4-31b & multi & s1 & self/no-base & 87.2 & 87.2 \\
gemma-4-31b & multi & s1 & neg/base-typ & 87.2 & 87.2 \\
gemma-4-31b & multi & s1 & neg/no-base & 87.2 & 87.2 \\
gemma-4-31b & multi & s2 & self/base-typ & 76.5 & 76.5 \\
gemma-4-31b & multi & s2 & self/no-base & 76.5 & 76.5 \\
gemma-4-31b & multi & s2 & neg/base-typ & 76.5 & 76.5 \\
gemma-4-31b & multi & s2 & neg/no-base & 76.5 & 76.5 \\
gemma-4-31b & multi & s3 & self/base-typ & 77.3 & 77.3 \\
gemma-4-31b & multi & s3 & self/no-base & 77.3 & 77.3 \\
gemma-4-31b & multi & s3 & neg/base-typ & 77.3 & 77.3 \\
gemma-4-31b & multi & s3 & neg/no-base & 77.3 & 77.3 \\
gemma-4-31b & multi & s4 & self/base-typ & 82.6 & 82.6 \\
gemma-4-31b & multi & s4 & self/no-base & 82.6 & 82.6 \\
gemma-4-31b & multi & s7 & neg/base-typ & 76.8 & 76.8 \\
gemma-4-31b & multi & s7 & neg/no-base & 76.8 & 76.8 \\
gemma-4-31b & multi & s13 & self/base-typ & 83.8 & 83.8 \\
gemma-4-31b & multi & s13 & self/no-base & 83.8 & 83.8 \\
gemma-4-31b & multi & s13 & neg/base-typ & 83.8 & 83.8 \\
gemma-4-31b & multi & s13 & neg/no-base & 83.8 & 83.8 \\
\midrule
qwen-3.5-9b & upper & s1 & self/base-typ & 78.6 & 78.6 \\
qwen-3.5-9b & upper & s1 & self/no-base & 78.6 & 78.6 \\
qwen-3.5-9b & upper & s1 & neg/base-typ & 78.6 & 78.6 \\
qwen-3.5-9b & upper & s1 & neg/no-base & 78.6 & 78.6 \\
qwen-3.5-9b & upper & s2 & self/base-typ & 74.3 & 74.3 \\
qwen-3.5-9b & upper & s2 & self/no-base & 74.3 & 74.3 \\
qwen-3.5-9b & upper & s2 & neg/base-typ & 74.3 & 74.3 \\
qwen-3.5-9b & upper & s2 & neg/no-base & 74.3 & 74.3 \\
qwen-3.5-9b & upper & s3 & self/base-typ & 78.5 & 78.5 \\
qwen-3.5-9b & upper & s3 & self/no-base & 78.5 & 78.5 \\
qwen-3.5-9b & upper & s3 & neg/base-typ & 78.5 & 78.5 \\
qwen-3.5-9b & upper & s3 & neg/no-base & 78.5 & 78.5 \\
qwen-3.5-9b & upper & s4 & self/base-typ & 76.4 & 76.4 \\
qwen-3.5-9b & upper & s4 & self/no-base & 76.5 & 76.5 \\
qwen-3.5-9b & upper & s7 & neg/base-typ & 78.3 & 78.3 \\
qwen-3.5-9b & upper & s7 & neg/no-base & 78.3 & 78.3 \\
qwen-3.5-9b & upper & s13 & self/base-typ & 76.4 & 76.4 \\
qwen-3.5-9b & upper & s13 & self/no-base & 76.5 & 76.5 \\
qwen-3.5-9b & upper & s13 & neg/base-typ & 76.5 & 76.5 \\
qwen-3.5-9b & upper & s13 & neg/no-base & 76.5 & 76.5 \\
\midrule
qwen-3.5-9b & multi & s1 & self/base-typ & 82.3 & 82.3 \\
qwen-3.5-9b & multi & s1 & self/no-base & 82.3 & 82.3 \\
qwen-3.5-9b & multi & s1 & neg/base-typ & 82.4 & 82.4 \\
qwen-3.5-9b & multi & s1 & neg/no-base & 82.3 & 82.3 \\
qwen-3.5-9b & multi & s2 & self/base-typ & 77.9 & 77.9 \\
qwen-3.5-9b & multi & s2 & self/no-base & 77.9 & 77.9 \\
qwen-3.5-9b & multi & s2 & neg/base-typ & 77.9 & 77.9 \\
qwen-3.5-9b & multi & s2 & neg/no-base & 77.9 & 77.9 \\
qwen-3.5-9b & multi & s3 & self/base-typ & 74.8 & 74.8 \\
qwen-3.5-9b & multi & s3 & self/no-base & 74.8 & 74.8 \\
qwen-3.5-9b & multi & s3 & neg/base-typ & 74.8 & 74.8 \\
qwen-3.5-9b & multi & s3 & neg/no-base & 74.8 & 74.8 \\
qwen-3.5-9b & multi & s4 & self/base-typ & 80.5 & 80.5 \\
qwen-3.5-9b & multi & s4 & self/no-base & 80.5 & 80.5 \\
qwen-3.5-9b & multi & s7 & neg/base-typ & 80.4 & 80.4 \\
qwen-3.5-9b & multi & s7 & neg/no-base & 80.4 & 80.4 \\
qwen-3.5-9b & multi & s13 & self/base-typ & 80.4 & 80.4 \\
qwen-3.5-9b & multi & s13 & self/no-base & 80.4 & 80.4 \\
qwen-3.5-9b & multi & s13 & neg/base-typ & 80.5 & 80.5 \\
qwen-3.5-9b & multi & s13 & neg/no-base & 80.4 & 80.4 \\
\midrule
qwen-3.5-9b (base) & upper & base & self/no-base & 65.3 & 65.3 \\
qwen-3.5-9b (base) & upper & base & neg/no-base & 65.3 & 65.3 \\
\midrule
qwen-3.5-9b (base) & multi & base & self/no-base & 66.8 & 66.8 \\
qwen-3.5-9b (base) & multi & base & neg/no-base & 67.0 & 67.0 \\
\bottomrule
\end{longtable}
}

{\small
\begin{longtable}{llll rr}
\caption{HumanEval TEST Pearson ($\times100$), raw vs tc gen variant. Systems: gemma-4-31b \& qwen-3.5-9b $\times$ base/s1/s2/s3/s4/s7/s13, both datasets.}\label{tab:he_test_pearson}\\
\toprule
Model & Data & Setting & Eval & raw & tc \\
\midrule
\endfirsthead
\multicolumn{6}{l}{\itshape (continued)}\\
\toprule
Model & Data & Setting & Eval & raw & tc \\
\midrule
\endhead
\bottomrule
\endfoot
gemma-4-31b & upper & s1 & self/base-typ & 51.1 & 31.8 \\
gemma-4-31b & upper & s1 & self/no-base & 51.1 & 31.4 \\
gemma-4-31b & upper & s1 & neg/base-typ & 51.1 & 42.2 \\
gemma-4-31b & upper & s1 & neg/no-base & 51.1 & 49.0 \\
gemma-4-31b & upper & s2 & self/base-typ & 59.2 & 41.0 \\
gemma-4-31b & upper & s2 & self/no-base & 59.2 & 41.2 \\
gemma-4-31b & upper & s2 & neg/base-typ & 59.2 & 55.5 \\
gemma-4-31b & upper & s2 & neg/no-base & 59.2 & 59.9 \\
gemma-4-31b & upper & s3 & self/base-typ & 67.4 & 50.5 \\
gemma-4-31b & upper & s3 & self/no-base & 67.4 & 50.5 \\
gemma-4-31b & upper & s3 & neg/base-typ & 67.4 & 60.2 \\
gemma-4-31b & upper & s3 & neg/no-base & 67.4 & 63.2 \\
gemma-4-31b & upper & s4 & self/base-typ & 60.8 & 63.3 \\
gemma-4-31b & upper & s4 & self/no-base & 60.8 & 52.5 \\
gemma-4-31b & upper & s7 & neg/base-typ & 58.8 & 75.8 \\
gemma-4-31b & upper & s7 & neg/no-base & 58.8 & 59.7 \\
gemma-4-31b & upper & s13 & self/base-typ & 53.2 & 37.8 \\
gemma-4-31b & upper & s13 & self/no-base & 53.2 & 35.8 \\
gemma-4-31b & upper & s13 & neg/base-typ & 53.2 & 49.9 \\
gemma-4-31b & upper & s13 & neg/no-base & 53.2 & 46.2 \\
\midrule
gemma-4-31b & multi & s1 & self/base-typ & 35.6 & 55.7 \\
gemma-4-31b & multi & s1 & self/no-base & 35.6 & 42.4 \\
gemma-4-31b & multi & s1 & neg/base-typ & 35.6 & 71.3 \\
gemma-4-31b & multi & s1 & neg/no-base & 35.6 & 43.0 \\
gemma-4-31b & multi & s2 & self/base-typ & 59.5 & 64.9 \\
gemma-4-31b & multi & s2 & self/no-base & 59.5 & 58.3 \\
gemma-4-31b & multi & s2 & neg/base-typ & 59.5 & 70.7 \\
gemma-4-31b & multi & s2 & neg/no-base & 59.5 & 53.2 \\
gemma-4-31b & multi & s3 & self/base-typ & 58.9 & 61.6 \\
gemma-4-31b & multi & s3 & self/no-base & 58.9 & 59.3 \\
gemma-4-31b & multi & s3 & neg/base-typ & 58.9 & 70.4 \\
gemma-4-31b & multi & s3 & neg/no-base & 58.9 & 58.1 \\
gemma-4-31b & multi & s4 & self/base-typ & 48.8 & 69.0 \\
gemma-4-31b & multi & s4 & self/no-base & 48.8 & 57.6 \\
gemma-4-31b & multi & s7 & neg/base-typ & 41.6 & 78.1 \\
gemma-4-31b & multi & s7 & neg/no-base & 41.6 & 57.6 \\
gemma-4-31b & multi & s13 & self/base-typ & 31.9 & 51.0 \\
gemma-4-31b & multi & s13 & self/no-base & 31.9 & 38.3 \\
gemma-4-31b & multi & s13 & neg/base-typ & 31.9 & 64.6 \\
gemma-4-31b & multi & s13 & neg/no-base & 31.9 & 39.6 \\
\midrule
qwen-3.5-9b & upper & s1 & self/base-typ & 41.5 & 45.4 \\
qwen-3.5-9b & upper & s1 & self/no-base & 41.5 & 40.9 \\
qwen-3.5-9b & upper & s1 & neg/base-typ & 41.5 & 56.3 \\
qwen-3.5-9b & upper & s1 & neg/no-base & 41.4 & 42.0 \\
qwen-3.5-9b & upper & s2 & self/base-typ & 65.8 & 55.3 \\
qwen-3.5-9b & upper & s2 & self/no-base & 65.8 & 46.0 \\
qwen-3.5-9b & upper & s2 & neg/base-typ & 65.8 & 69.5 \\
qwen-3.5-9b & upper & s2 & neg/no-base & 65.8 & 68.9 \\
qwen-3.5-9b & upper & s3 & self/base-typ & 52.6 & 49.9 \\
qwen-3.5-9b & upper & s3 & self/no-base & 52.7 & 43.2 \\
qwen-3.5-9b & upper & s3 & neg/base-typ & 52.7 & 67.6 \\
qwen-3.5-9b & upper & s3 & neg/no-base & 52.7 & 65.5 \\
qwen-3.5-9b & upper & s4 & self/base-typ & 46.4 & 58.9 \\
qwen-3.5-9b & upper & s4 & self/no-base & 46.4 & 51.5 \\
qwen-3.5-9b & upper & s7 & neg/base-typ & 41.4 & 68.5 \\
qwen-3.5-9b & upper & s7 & neg/no-base & 41.4 & 55.8 \\
qwen-3.5-9b & upper & s13 & self/base-typ & 43.9 & 41.7 \\
qwen-3.5-9b & upper & s13 & self/no-base & 43.9 & 37.1 \\
qwen-3.5-9b & upper & s13 & neg/base-typ & 43.9 & 53.5 \\
qwen-3.5-9b & upper & s13 & neg/no-base & 43.9 & 38.2 \\
\midrule
qwen-3.5-9b & multi & s1 & self/base-typ & 25.9 & 57.8 \\
qwen-3.5-9b & multi & s1 & self/no-base & 25.9 & 48.4 \\
qwen-3.5-9b & multi & s1 & neg/base-typ & 25.9 & 65.8 \\
qwen-3.5-9b & multi & s1 & neg/no-base & 25.9 & 35.1 \\
qwen-3.5-9b & multi & s2 & self/base-typ & 67.2 & 64.7 \\
qwen-3.5-9b & multi & s2 & self/no-base & 67.2 & 61.3 \\
qwen-3.5-9b & multi & s2 & neg/base-typ & 67.2 & 70.1 \\
qwen-3.5-9b & multi & s2 & neg/no-base & 67.2 & 75.2 \\
qwen-3.5-9b & multi & s3 & self/base-typ & 55.0 & 64.6 \\
qwen-3.5-9b & multi & s3 & self/no-base & 55.0 & 57.8 \\
qwen-3.5-9b & multi & s3 & neg/base-typ & 55.0 & 72.4 \\
qwen-3.5-9b & multi & s3 & neg/no-base & 55.0 & 61.8 \\
qwen-3.5-9b & multi & s4 & self/base-typ & 32.9 & 66.0 \\
qwen-3.5-9b & multi & s4 & self/no-base & 32.9 & 56.5 \\
qwen-3.5-9b & multi & s7 & neg/base-typ & 15.6 & 67.9 \\
qwen-3.5-9b & multi & s7 & neg/no-base & 15.6 & 40.3 \\
qwen-3.5-9b & multi & s13 & self/base-typ & 27.5 & 53.1 \\
qwen-3.5-9b & multi & s13 & self/no-base & 27.5 & 44.4 \\
qwen-3.5-9b & multi & s13 & neg/base-typ & 27.5 & 63.6 \\
qwen-3.5-9b & multi & s13 & neg/no-base & 27.5 & 21.2 \\
\midrule
qwen-3.5-9b (base) & upper & base & self/no-base & 32.1 & 26.5 \\
qwen-3.5-9b (base) & upper & base & neg/no-base & 32.1 & 30.4 \\
\midrule
qwen-3.5-9b (base) & multi & base & self/no-base & 4.5 & 27.1 \\
qwen-3.5-9b (base) & multi & base & neg/no-base & 4.4 & 29.3 \\
\bottomrule
\end{longtable}
}


\end{document}
